\section{Introduction}
\label{sec:introduction}

% Spiking Neural Networks (SNNs) transmit information through discrete spikes over time, enabling efficient, event-driven processing~\cite{cao:2015, esser:2016, neftci:2018, roy:2019, cramer:2022, rathi:2023}, especially when supported by neuromorphic hardware~\cite{merolla:2014, akopyan:2015, davies:2018}. Recent advances in direct training methods~\cite{wu:2018, wu:2019} have significantly improved the accuracy of SNNs, even with a limited number of timesteps, further enhancing their efficiency. However, despite these advancements, SNNs still exhibit a noticeable accuracy gap compared to deep neural networks. Moreover, the complex non-linear dynamics of spike-based processing can lead to degraded performance on large-scale and high-complexity datasets~\cite{review_guo:2023, review_zhou:2024}. Therefore, continued research is essential to enhance reliability and accuracy, ultimately paving the way for broader adoption in real-world applications.

Spiking Neural Networks (SNNs) transmit information through discrete spikes over time, enabling efficient, event-driven processing~\cite{cao:2015, esser:2016, neftci:2018, roy:2019, cramer:2022, rathi:2023}, especially on neuromorphic hardware~\cite{merolla:2014, akopyan:2015, davies:2018}. Recent advances in direct training~\cite{wu:2018, wu:2019} have significantly improved SNN accuracy with few timesteps, enhancing efficiency. Yet, SNNs still lag behind deep neural networks, and their complex dynamics often hinder performance on large-scale datasets~\cite{review_guo:2023, review_zhou:2024}. Thus, further innovations are crucial to ensure reliability and accuracy for real-world adoption.

One of the primary challenges in SNN training is \textbf{Temporal Covariate Shift (TCS)}~\cite{bntt:2020}, where internal covariate shifts in activations accumulate over time. While TEBN~\cite{tebn:2022} and TAB~\cite{tab:2024} attempt to address this issue by introducing timestep-specific normalization parameters, they disrupt parameter homogeneity across time, increasing both model complexity and computational cost. More importantly, our analysis reveals that these methods fail to mitigate TCS in other critical components of SNNs, particularly the membrane potential. Since the membrane potential directly determines neuronal output, neglecting TCS in this component can lead to significant accuracy degradation.

The second challenge lies in \textbf{training the internal parameters}~\cite{fang:2021, wang:2022, rathi:2023} of Leaky Integrate-and-Fire (LIF) neurons, which are the threshold voltage ($V_{\text{thr}}$) and leakage term ($\tau$). Despite their crucial role in LIF neuron behavior, they are typically fixed or updated via gradient descent under strict constraints. We reveal that existing gradient flow heavily depends on the scale of $V_{\text{thr}}$, and unconstrained updates during training could lead to significant instability. Surprisingly, none of the existing studies has explicitly recognized this issue.

In this work, we introduce a comprehensive approach to overcoming these limitations. Our first contribution, \textbf{MP-Init (Membrane Potential Initialization)}, stems from our analysis of the root cause of TCS in SNNs. By addressing TCS in membrane dynamics, MP-Init ensures consistent neuron behavior across timesteps without using timestep-dependent parameters. Our second contribution, \textbf{TrSG (Threshold-Robust Surrogate Gradient)}, is a refined version of surrogate gradient that stabilizes the gradient flow with respect to the threshold $V_{\text{thr}}$. This modification allows for stable gradient-based updates of all internal parameters, leading to superior convergence.

We provide a rigorous mathematical analysis of our methods and demonstrate through extensive experiments that they significantly improve SNN performance. Our approach enables effective SNN training, achieving state-of-the-art results on both static and dynamic datasets.